\section{Introduction}
\label{sec:intro}

The marginal value of public funds (MVPF) has become a common currency for comparing government policies. It divides the beneficiaries' willingness to pay for a policy by the policy's net cost to the government, where the net cost includes the fiscal externalities created by the behavioral responses the policy induces \citep{hendren2016policy,hendren2020unified}. Because it puts a tax cut, a health insurance subsidy, and a training program on the same scale, the MVPF identifies which policies deliver the most welfare per dollar and which pay for themselves. Most existing MVPFs rest on quasi-experimental estimates of behavioral responses, and their long-run fiscal externalities are projected rather than observed \citep{hendren2020unified,Finkelstein2020welfare}.

Life-cycle policies fit this approach less easily, and retirement policy is the leading example. Social Security benefits depend on a worker's highest 35 years of earnings, so a change in the benefit formula changes the return to work at age 35 as well as at age 65. Rules that tie benefits to the claiming age change when people stop working, and therefore change income and payroll tax revenue and the outlays of disability insurance and means-tested transfers. Who responds depends on health: whether a worker can postpone retirement depends on how costly work becomes as health deteriorates. A natural experiment that shifts benefits for one cohort identifies a local response of claiming or employment. It does not identify the present value of the fiscal externalities that accumulate over a lifetime, nor the insurance value of a reform to a worker who does not yet know whether their health will hold up.

Retirement policy also carries large fiscal stakes. Population aging puts pressure on the financing of public pensions \citep{trustees2024,cbo2024socialsecurity}, and the most common response is to raise the age at which full benefits are paid. The Social Security Amendments of 1983 raised the normal retirement age (NRA) in the United States from 65 to 67, beginning with workers born in 1938. They also raised the delayed retirement credit (DRC), the increase in annual benefits for each year a worker postpones claiming beyond the NRA, gradually to 8 percent per year \citep{Svahn1983social,diamond1999social}. Cross-country simulations of pension reform likewise contrast increases in eligibility ages with changes in the actuarial adjustment of benefits \citep{gruber2009social}. This paper asks which retirement reforms deliver the most welfare per dollar of public funds, and which pay for themselves.

We answer this question in an estimated life-cycle model of retirement. Individuals face health shocks with three states---good health, poor health, and disability---that lower wages, raise out-of-pocket medical expenses, and raise mortality. They face persistent wage risk, choose consumption, saving, labor force participation, and hours, and decide when to claim Social Security benefits. Their preferences over consumption and leisure are nonseparable, as in \citet{french2011effects}, and working entails a fixed time cost that is larger in poor health, so that labor supply is nonconvex and retirement emerges as a choice. Reentering the labor market after an exit is costly. The model incorporates the main features of the US system: the Social Security benefit formula with its claiming-age adjustments and earnings test, disability insurance, employer pensions, Medicare, a progressive and age-dependent income tax, payroll taxes, and a means-tested consumption floor. We estimate it in two steps on the Panel Study of Income Dynamics and the Medical Expenditure Panel Survey for the cohorts born between 1920 and 1940. Health transitions, mortality, medical expenses, and wages are estimated directly from the data, and preferences are estimated by the method of simulated moments to match life-cycle profiles of participation and hours by health, of mean and median assets, and of labor market reentry. The estimated model reproduces the main features of these profiles: high and stable participation in the thirties and forties, an earlier decline of participation among workers in poor health or with a disability beginning in their fifties, and the sharp exits of the sixties.

The estimates imply that working requires a fixed time cost of 791 hours per year in good health and 868 hours in poor health or with a disability---ratios of about 18 and 19 percent to the estimated annual time endowment of 4,491 hours---and that returning to work after an exit costs 3,312 hours. These nonconvexities make retirement a discrete and, in practice, irreversible choice, and they determine which workers can respond when the reward for working an additional year changes. In our companion paper \citep{NdiayeYu2026frot}, we show that these nonconvexities are the main driver of nonparticipation at older ages in this economy, ahead of the effects of health on wages and on mortality.

We use the estimated model to compute the MVPF of 66 reforms in five families: the claiming-age dependence of Social Security benefits (the NRA, the DRC, and symmetric adjustments around the NRA), the progressivity and level of the benefit formula, Medicare coverage after age 65, income taxes after age 62, and the return to saving. For each reform, we compute the ex ante willingness to pay of a worker at age 30, before health and wage histories are realized, the change in Social Security's own outlays, the change in the net cost of other taxes and transfers, and the implied MVPF. The model thus serves as a laboratory in which every component of the MVPF of a life-cycle policy can be computed, including fiscal externalities that materialize decades after behavior changes and the insurance value of a reform against health risk.

Our central result is that, for the cohorts we study, the MVPF of a claiming-age reform depends on whether the reform rewards later claiming or lowers benefits at every claiming age. Consider first the reforms that reward later claiming. Raising the DRC from the 4.5 percent faced by these cohorts to between 8 and 10 percent, or making the adjustment of benefits around the NRA symmetric and steeper, raises welfare by 0.10 to 0.80 percent of consumption at every age and lowers net government spending by 0.12 to 1.57 percent. These reforms have an infinite MVPF: they pay for themselves. Raising the NRA to 66, 67, or 68 instead lowers welfare by 0.38 to 1.09 percent while saving 0.87 to 2.38 percent of government spending. Its MVPF, 0.16 to 0.18, is the welfare cost of each dollar saved, and it is lower than the MVPF of every benefit expansion in our grid except the Medicare expansions. Raising the NRA is therefore a comparatively inexpensive way to save money, but rewards for later claiming save money while raising welfare. We call an NRA increase a mandate in a loose sense: it does not prevent anyone from claiming at 62, but it lowers the benefit at every claiming age, so that a worker must claim later to receive a given benefit.

The welfare and fiscal sides of this result have different explanations. On the welfare side, a higher DRC weakly enlarges every worker's choice set at given prices and taxes: it raises the benefit of anyone who claims after the NRA and lowers no one's. Its value, assessed at age 30 before health is known, combines a higher reward for working longer when work remains feasible with unchanged terms for workers whose health leads them to claim early. A higher NRA lowers the benefits of all claimants, including those who cannot postpone retirement at low cost; the additional participation it generates at ages 65 to 69, 0.60 to 3.27 percentage points, is a response to lower benefits rather than to a reward. The fiscal side is the substantive result. Social Security's own outlays fall under every DRC increase, by 2.46 to 5.05 percent, even though a higher credit would raise outlays at unchanged behavior; the claiming responses therefore more than finance the higher credit within the program. A one-year delay lowers the present value of a claimant's benefits whenever the credit is below the rate that would make the delay actuarially fair from the government's perspective, given the claimant's survival prospects and the interest rate, and this rate is higher for claimants in worse health. The fall in outlays implies that the induced delays occur, on balance, where the credit is below that rate, in line with evidence that delaying beyond the NRA was not actuarially advantageous for cohorts that faced credits well below 8 percent \citep{Coile2002delays,Shoven2014does}. The moderate versions of these reforms also raise participation in every age group we track, and the net cost of the other taxes and transfers falls as careers lengthen.

Two of our reforms correspond to the end points of the two changes to claiming incentives in the 1983 Amendments: an NRA of 67 and a DRC of 8 percent. Evaluated for the 1920--1940 cohorts, whose schedule we represent with an NRA of 65 and a cohort-average credit of 4.5 percent, the credit increase raises welfare by 0.16 percent and saves 0.45 percent of government spending, whereas the increase in the NRA lowers welfare by 0.74 percent and saves 1.75 percent, an MVPF of 0.16. The two halves of the change sit at opposite ends of the ranking. In our grid, rewarding later claiming also dominates raising the NRA whenever the required fiscal saving is moderate: raising the DRC to 10 percent saves nearly as much as raising the NRA to 66 while raising welfare, and the steepest symmetric adjustment we consider saves more than the increase to 66 and raises welfare by 0.80 percent.

The remaining reform families sharpen the picture. Every expansion of benefit generosity except the smallest raises net spending. The smallest increase in the bottom replacement rate of the benefit formula is approximately budget neutral, because the taxes paid by the additional workers it keeps in the labor force offset its direct cost; larger increases in the bottom rate have MVPFs between 0.69 and 1.14, above those of proportional increases in all rates, between 0.30 and 0.61. The bottom rate targets workers whose lifetime earnings were depressed by low ability or poor health. Raising the upper replacement rates on top of the bottom one adds little: from the largest increase in the bottom rate to the largest increase in the second rate, willingness to pay rises by 0.15 percent of consumption while the net cost rises by 0.83 percent of government spending. Expanding Medicare coverage after age 65 by 20 percentage points has the lowest MVPF of any expansion, about 0.15. Cutting income taxes after age 62 generates the largest participation responses of any reform we study, up to 8.9 percentage points at ages 65 to 69, yet its MVPF is only 0.23 to 0.26. Subsidies to saving raise welfare at a small net cost, because additional saving raises tax revenue, and they lower participation at older ages; we show that their MVPFs rest on a different consumption base from the other reforms. Finally, eliminating each instrument in turn shows that Social Security is the most valuable of them: without it, welfare would be 5.89 percent of consumption lower. Because such eliminations do not hold the budget fixed, we rely on the MVPF to compare reforms.

These results suggest an order of operations for closing a financing gap under the rules we study. The MVPF ranks the ways of saving a dollar by their welfare cost. Reforms that reward later claiming come first: they save money and raise welfare at the same time. To a first-order approximation that reads our simulated expansions in reverse, raising the NRA costs less welfare per dollar saved than raising income taxes after age 62 or cutting all replacement rates proportionally, but it is dominated by rewards for later claiming up to the largest fiscal saving that a reward-delay reform in our grid achieves.

Three features of the MVPF of a life-cycle policy drive these conclusions and are easy to miss in short-run evaluations. First, willingness to pay is computed ex ante, before health is known, so it values the protection a reform provides, or removes, for workers whose health later deteriorates. Second, a substantial part of the fiscal consequences of retirement reforms falls outside Social Security and accrues over many years after the reform first changes behavior: the smallest expansion of the bottom replacement rate is budget neutral only because of the taxes paid by additional workers. Third, MVPF levels depend on how the willingness to pay is converted into dollars. We value it over consumption from age 62 onward for reforms that operate through retirement benefits, Medicare, or old-age taxes. This convention leaves every ranking among these reforms and every infinite MVPF unchanged, but the finite levels are about one-fifth of those obtained by valuing the willingness to pay over lifetime consumption, the normalization under which a transfer with no behavioral response has an MVPF of one (Section~\ref{sec:discussion}). We report every component of the MVPF separately and state the formula that maps them into it; every reported MVPF equals the reported ratio of willingness to pay to net cost times a common factor that the tables pin down.

\paragraph{Related literature.} This paper contributes to four literatures. First, it builds on structural life-cycle models of retirement, which have shown how Social Security, pensions, health, and health insurance shape labor supply and saving at older ages \citep{Gustman1986structural,rust1989dynamic,Rust1997how,french2005effects,vanderKlaauw2008social,Denardi2010,french2011effects,Haan2014,Low2015,blundell2016retirement,Denardi2018,fan2024estimation}, and which have been used to evaluate changes in the NRA and other Social Security reforms \citep{Conesa1999social,Imrohoroglu2012,Kitao2014sustainable,Laun2015life,Li2018,Bairoliya2023}. \citet{Jones2023} find that, with mortality that differs by education, the welfare-maximizing reform of the current system weakens the dependence of benefits on lifetime earnings and reduces claiming-age adjustments. Their exercise searches for the best reform of the current system; ours measures the value per dollar of marginal reforms from the cohort-average credit of 4.5 percent, and we discuss the relation between the two in Section~\ref{sec:discussion}. Our model follows \citet{french2005effects} and \citet{french2011effects} in its treatment of preferences, wages, and the Social Security rules. It adds a disability state with disability insurance, which provides an alternative pathway out of the labor force at older ages \citep{Maestas2013,French2014,Low2015,Yu2024}, fixed costs of work that depend on health, and age-dependent progressive income taxes \citep{Borella2023}, and it targets the population rather than male household heads. Relative to this literature, our contribution is to map a large set of reforms into a common welfare metric that accounts for their fiscal consequences across programs.

Second, the paper contributes to the welfare analysis of government policies with the MVPF \citep{hendren2016policy,hendren2020unified} and to the sufficient-statistics approach to social insurance \citep{Chetty2009sufficient,Kleven2021sufficient}. Closest to our question, \citet{haller2022welfare} evaluates pension reforms with sufficient statistics in the Baily--Chetty tradition, where the welfare effect of a reform compares a consumption-smoothing benefit with the total fiscal effect of the reform, and \citet{Kolsrud2024retirement} show with Swedish data that workers who retire earlier consume less in retirement, so that reforms pushing workers to retire later carry a substantial consumption-smoothing cost. We show how the components of our MVPF map into the Baily--Chetty comparison. The structural approach measures the fiscal externality over the life cycle and across programs and computes the ex ante insurance value of each reform; sufficient-statistics estimates discipline the local behavioral responses on which such calculations rest. The two approaches are complements.

Third, the paper relates to the evidence on how workers respond to Social Security rules. Program rules generate strong early-retirement incentives across countries \citep{gruber1998social,Gruberwise2002}, and older Americans have been working longer \citep{Coile2019}. Workers respond to benefit levels, to the marginal return to additional work, and to the early and normal retirement ages \citep{Coile2007future,Liebman2009labor,Mastrobuoni2009labor,Staubli2013does,Manoli2016nonparametric}. Claiming and labor force exit are closely linked \citep{shoven2017}; an information intervention raised labor force participation without detectable effects on claiming \citep{liebman2015would}; the earnings test affects employment \citep{gelber2020employment,Gelber2020estimating}; and an increase in the NRA raises disability insurance enrollment \citep{Duggan2007aching}. Statutory retirement ages also affect behavior beyond the financial incentives they create: they act as reference points for retirement \citep{behagel2012,seibold2021reference,Lalive2023how}, and when the NRA rose from 65 to 66, claiming moved with it while labor force exit remained concentrated at 65 \citep{Deshpande2024how}. Our reforms of the claiming-age schedule work only through financial incentives; we discuss how reference-point effects of statutory ages would change our conclusions.

Fourth, the paper complements the normative analysis of retirement and social insurance \citep{diamond1978model,cremer2004social,golosov2006designing,huggett2010well,Golosov2013optimal,hosseini2015adverse,nishiyama2007does,hosseini2019retirement}. In our companion paper \citep{NdiayeYu2026frot}, we characterize the constrained-efficient allocation in a dynamic Mirrlees economy with endogenous retirement, estimated on the same data, and find that optimal retirement consumption rises with the claiming age more steeply than the statutory delayed retirement credit implies. The present paper takes the opposite route. Instead of characterizing the optimum, it evaluates reforms within the existing system. The two analyses point in the same direction: the reforms that steepen the claiming-age schedule are those that pay for themselves.

\paragraph{Outline.} Section~\ref{sec:model} presents the model. Section~\ref{sec:estimation} describes the data and the estimation and evaluates the model's fit. Section~\ref{sec:framework} defines the welfare and fiscal measures, the MVPF, and its decomposition. Section~\ref{sec:results} presents the results, and Section~\ref{sec:financing} draws their implications for closing a financing gap. Section~\ref{sec:discussion} discusses the measurement conventions, general equilibrium, reference-point effects of statutory ages, and the relation between the cohorts we study and current law. Section~\ref{sec:conclusion} concludes. The appendices contain the details of the model, the data, the estimation, the construction of the reforms, and the derivations.
